% 土星（综述）
% license CCBYSA3
% type Wiki

本文根据 CC-BY-SA 协议转载翻译自维基百科\href{https://en.wikipedia.org/wiki/Saturn}{相关文章}。

\begin{figure}[ht]
\centering
\includegraphics[width=8cm]{./figures/13ddb48d3e7de073.png}
\caption{土星及其显著的行星环，由卡西尼轨道器拍摄\(^\text{[a]}\)} \label{fig_Saturn_1}
\end{figure}
土星是距离太阳第六颗行星，也是太阳系中仅次于木星的第二大行星。它是一颗气态巨行星，其平均半径约为地球的 9 倍。土星的平均密度仅为地球的八分之一，但质量超过地球的 95 倍。尽管土星几乎与木星一样大，但其质量还不到木星的三分之一。土星以 9.59 AU（14.34 亿公里）的距离绕太阳运行，其轨道周期为 29.45 年。

土星内部被认为由一颗岩石核心组成，外围为一层深厚的金属氢层，其上是液态氢与液态氦的中间层，以及最外层的气体外壳。土星呈现淡黄色调，是由于其上层大气中的氨晶体所致。金属氢层中的电流被认为会产生土星的行星磁场，该磁场强度弱于地球，但由于土星体积更大，其磁矩为地球的 580 倍。土星磁场强度约为木星的二十分之一。土星外层大气整体上较为平淡，缺乏对比度，但也可能出现长期存在的大气结构。土星的风速可达每小时 1,800 公里（每小时 1,100 英里）。

这颗行星拥有明亮且庞大的环系统，主要由冰粒组成，且掺杂少量岩石碎片与尘埃。至少有 274 颗卫星环绕土星运行，其中 63 颗已被正式命名；这一数字不包括土星光环中的数百颗小卫星体（moonlets）。土卫六（Titan）是土星最大的卫星，也是太阳系中第二大的卫星，其体积大于（但质量小于）水星，并且是太阳系中唯一拥有稠密大气层的卫星。

名称与符号

土星以罗马财富与农业之神之名命名，而该神是朱庇特之父。其天文符号（♄）可追溯到希腊的《奥克绪林库斯纸草文稿》（Oxyrhynchus Papyri），其中该符号可被视为希腊字母 κ–ρ（kappa–rho）的连写，并带有一条水平短线，用作 Κρονος（Cronus，即土星的希腊名称）的缩写。后来该符号演变为类似小写希腊字母 η（eta）的形式，并在 16 世纪时于顶部加上十字，以将这一异教符号基督化。

罗马人以行星土星之名为一周的第七天命名为 Saturday，即 Sāturni diēs，“土星之日”。

\begin{figure}[ht]
\centering
\includegraphics[width=8cm]{./figures/91b99dbcbbce021d.png}
\caption{} \label{fig_Saturn_2}
\end{figure}



